\section{Introduction}
\label{sec:introduction}

The Heavy Photon Search (HPS) experiment at the Thomas Jefferson National Accelerator
Facility (Jefferson Lab) searches for light, weakly coupled particles produced in
electron fixed-target scattering and decaying to charged particles in a compact
forward spectrometer. A benchmark target is a visibly decaying dark photon, denoted
$\Ap$, that kinetically mixes with the Standard Model photon with coupling strength
$\eps$ \cite{Holdom1986,EssigEtAl2011,Fabbrichesi2020DarkPhoton}. In the minimal
visible model, an $\Ap$ lighter than the dimuon threshold decays promptly to
$e^+e^-$ for sufficiently large $\eps^2$, producing a narrow resonance in the
reconstructed electron--positron invariant-mass spectrum. Its natural width is much
smaller than the HPS mass resolution, so the observed line shape is set by the
detector response.

The prompt HPS search is built on a simple physical normalization. In fixed-target
electroproduction, the dark-photon yield is proportional to the radiative trident
rate,
\begin{equation}
\frac{dN_{\Ap}}{dm}\bigg|_{m=m_{\Ap}}
=
\frac{3\pi}{2\alpha_{\mathrm{EM}}}
m_{\Ap}\eps^2
\frac{dN_{\gamma^\ast}}{dm}\bigg|_{m=m_{\Ap}},
\label{eq:intro-ap-yield}
\end{equation}
below the dimuon threshold. The experimental task is therefore to extract or limit a
narrow signal yield on top of the smooth quantum electrodynamics (QED) background, and
then convert that yield to $\eps^2$.

HPS has already published prompt resonance searches using the 2015 and 2016
engineering-run data sets. The 2015 search used 1170 nb$^{-1}$ at a beam energy of
1.056 GeV and scanned 19--81 MeV, finding no significant excess
\cite{HPS2015DarkPhoton}. The 2016 search used 10608 nb$^{-1}$ at 2.3 GeV and scanned
39--179 MeV, again finding no significant prompt resonance
\cite{HPS2016PromptLong}. The broader HPS program also includes displaced-vertex dark
sector searches, including the recent search for dark-sector strongly interacting
massive particles (SIMPs) in the 2016 data set \cite{HPSSIMP2025}. The present work
keeps the same HPS publication style and detector language, but addresses a different
analysis problem: a prompt visible dark-photon search with a modern probabilistic
background model.

The main methodological advance is the use of Gaussian process regression (GPR) for
the smooth trident background. For each tested mass, bins near the signal hypothesis
are excluded from training. The Gaussian process is conditioned on the sideband bins
and predicts the background mean and covariance in the excluded window. This replaces
a local analytic background function with a probability distribution over smooth
backgrounds. The resulting covariance is not discarded: it enters the profiled
likelihood as correlated nuisance parameters, so the upper limits and local
significances know how uncertain the sideband interpolation is.

Figure~\ref{fig:intro-context} shows the visible-dark-photon parameter space used as
context for the HPS prompt program. The figure is a reference and comparison plot, not
the defining inference for the analysis. The inference is performed directly in the
selected HPS mass spectra using the likelihood described in Sec.~\ref{sec:gpr-stats}.

\begin{figure*}[t]
\centering
\ifigure{0.82\textwidth}{hps_published_context.png}
\caption{Prompt visible-dark-photon context for the HPS engineering-run results. The
highlighted HPS curves show the published 2015 and 2016 prompt limits, displayed in
the same low-mass phase-space convention used by the analysis note. Existing
constraints and thermal-relic targets are shown for orientation
\cite{HPS2015DarkPhoton,HPS2016PromptLong,BaBarDarkPhoton2014,KLOE2013,KLOE22015,
A1MAMI2014,APEX2011,NA48DarkPhoton2015,HADESDarkPhoton2014,PHENIXDarkPhoton2015,
NA64DarkPhoton2018,NA64X17Visible2021,LHCbDarkPhoton2020,DarkCast2018}.}
\label{fig:intro-context}
\end{figure*}

The paper is organized as follows. Section~\ref{sec:detector-data} summarizes the HPS
detector configurations and data samples. Section~\ref{sec:selection-normalization}
describes the prompt event selection, mass resolution, and radiative normalization.
Section~\ref{sec:gpr-stats} gives the GPR background model and statistical inference.
Section~\ref{sec:validation} summarizes the validation program. Section~\ref{sec:results}
presents the current staged limits and the comparison figures that will become the
publication result once the final inputs are frozen.
